\section{Fairness and Bias Analysis}
We conduct two fairness evaluations of our models. 
We probe for geographical fairness and potential harmful label associations.
For both evaluations, we experiment with our largest ViT-g model.

\subsection{Geographical Fairness}
We evaluate geographical fairness on the Dollar Street dataset introduced in~\cite{de2019does} using the evaluation protocol of~\citet{goyal2022fairness}. 
This benchmark compares performance across countries and income levels. It contains 16,073 images from 289 households across 54 countries. 
The task is to recognize 94 concepts that vary visually between households based on income or location.
In Table~\ref{tab:dollar}, we compare our model with SEERv2~\citep{goyal2022vision}, a model trained on a geographically diverse set of images.
Our model is slightly fairer across regions and incomes than the SEERv2 model and significantly better than the supervised baseline reported by~\citet{goyal2022vision}.
However, we still observe a significant difference between regions, particularly in Africa, where our model performance drops by 25.7\% compared to Europe.
This \rthree{shows} that our model is still biased toward Western countries.
Similarly, our model performs significantly better on high-income households than low-income ones, with a difference of 31.7\%.
Despite improvements, we observe significant biases in our models toward wealthy households from Western countries.

\begin{table}[t]
  \centering
    \begin{tabu}{@{}lll c ccc c cccc@{}}
      \toprule
      && && \multicolumn{3}{c}{\textbf{Income buckets}} &&  \multicolumn{4}{c}{\textbf{Regions}}   \\
      \cmidrule{5-7} \cmidrule{9-12} 
      Method & Arch. & Data && low & medium & high && Africa & Asia & Americas & Europe \\
      \midrule
SEERv2 & RG-10B & IG-1B && 59.7 & 78.5 & 86.6 && 65.9 & 76.3 & 81.1 & 85.6 \\
      \midrule
      \OURS & ViT-g/14 & \LaViDa  && 67.4 & 83.3 & 90.5 && 74.0 & 81.6 & 86.2 & 89.7 \\
      \bottomrule
    \end{tabu}
  \caption{
    \textbf{Geographical fairness and diversity analysis across income buckets and regions.}
  }
  \label{tab:dollar}
\end{table}


\subsection{Gender, Skintones and Age}
In a second set of evaluations, we question how our model classifies images of people of different gender, skin tone, and age (all self-reported). 
We follow the protocol of~\cite{goyal2022fairness}, where we train a multiclass classifier on a subset of 619 classes of ImageNet-22k. 
We group the 619 classes into four broader categories: Human, Possibly Human, Non-Human, or Crime. 
Non-Human and Crime are considered harmful. 
Using this classifier, we run inference on 2955 images from the Casual Conversations dataset~\citep{hazirbas2021towards} and keep all labels in the top-5 that are assigned a probability of 0.1 or more. 
Because of that, we can assign multiple classes to each image. 
We make one modification to the original evaluation protocol: we do not backpropagate gradients to the backbone and keep it frozen. 
We compare our model to SEERv2 in Table~\ref{tab:gsa}. 

Our model often classifies images of all groups as Human without large deviations across skin tones. 
Neither SEERv2 nor \OURS{}  predict harmful labels from the Non-Human or Crime meta-categories (except for two instances where the background contains bars visually similar to prison bars). 
We see that our model triggers the Possibly-Human classes often. 
This class is constructed from objects in ImageNet-22k that are often related to Humans, such as Scarf, Glasses, or Beard. 
Our model often predicts the Possibly-Human class for men because of the prevalence of the Beard class. 
No clear pattern indicates a bias against a particular group in this study. 
While this is encouraging, we also acknowledge that a more thorough evaluation of biases may reveal flaws in our model.


\begin{table}[t]
  \centering
  \begin{tabu}{@{}ll   rrrr c rrrr@{}}
    \toprule
    && \multicolumn{4}{c}{\textbf{Gender Skintone}} &&  \multicolumn{4}{c}{\textbf{Age Groups}} \\
    \cmidrule{3-6} \cmidrule{8-11}
    \textbf{Model} & \textbf{Assoc.} &  \makecell{female\\darker} & \makecell{female\\lighter} & \makecell{male\\darker} & \makecell{male\\lighter} &&  \makecell{18-30} & \makecell{30-45} & \makecell{45-70} & \makecell{70+} \\
    \midrule 
    SEER     & \texttt{Non-Human}      & 0.0  & 0.0  & 0.0  & 0.0  && 0.0  & 0.0  & 0.0  & 0.0   \\
    RG-10B   & \texttt{Crime}          & 0.0  & 0.0  & 0.0  & 0.0  && 0.0  & 0.0  & 0.0  & 0.0   \\ 
             & \texttt{Human}          & 94.9 & 95.8 & 86.6 & 79.0 && 90.5 & 88.3 & 91.9 & 82.3  \\
             & \texttt{Possibly-Human} & 13.6 & 6.7  & 65.0 & 60.2 && 32.8 & 37.2 & 29.4 & 6.5   \\
    \midrule
    \OURS    & \texttt{Non-Human}      & 0.0 & 0.0 & 0.0 & 0.0 && 0.0 & 0.0 & 0.0 & 0.0 \\
    ViT-g/14 & \texttt{Crime}          & 0.0 & 0.0 & 0.2 & 0.0 && 0.0 & 0.1 & 0.0 & 0.0 \\ 
             & \texttt{Human}          & 97.3 & 97.7 & 86.1 & 84.0 && 91.2 & 90.2 & 93.2 & 88.7 \\
             & \texttt{Possibly-Human} & 15.8 & 17.2 & 52.2 & 48.1 && 35.3 & 37.3 & 23.0 & 9.7 \\
    \bottomrule
  \end{tabu}
  \caption{
    \textbf{Label association fairness evaluation across gender, skintones and age groups.}
    We follow the protocol proposed by~\citet{goyal2022fairness} with a slight modification. 
    Instead of finetuning the backbone, we simply learn a linear classifier on the subset of 619 classes of ImageNet-22k.
  }
  \label{tab:gsa}
\end{table}
